\subsection{半单李代数中的半单元与幂零元}

命题 3. 设 M 是 K 上的有限维向量空间，g 是 gl(M) 的半单子代数。则 g 包含其元素的半单分量和幂零分量。

若 \(\mathbf{K}_1\) 是 \(\mathbf{K}\) 的扩域，则 \(\mathfrak{g}_{(\mathbf{K}_1)}\) 的 Killing 形式是向 \(\mathfrak{g}_{(\mathbf{K}_1)}\) 延拓的 \(\mathfrak{g}\) 的 Killing 形式（\(\S 3\)，第 8 号），因而是非退化的；所以 \(\mathfrak{g}_{(\mathbf{K}_1)}\) 是半单的。因此只需在基域代数闭时证明命题 3，以下均假设情形如此。

对 \(\mathbf{N}\) 的每个子空间 \(\mathbf{M}\)，令 \(\mathfrak{g}_{\mathbf{N}}\) 为 \(\mathfrak{gl}(\mathbf{M})\) 中由使 \(\mathbf{N}\) 保持稳定且在 \(\mathbf{N}\) 上的限制迹为零的元素组成的子代数。由于 \(\mathfrak{g} = \mathcal{D}\mathfrak{g}\)，\(\mathfrak{g} \subset \mathfrak{g}_{\mathbf{N}}\)；当 \(\mathbf{N}\) 在 \(\mathfrak{g}\) 下稳定时，后者成立。再令 \(\mathfrak{g}^*\) 为 \(\mathfrak{g}\) 在 \(\mathfrak{gl}(\mathbf{M})\) 中的规范化子与各个 \(\mathfrak{g}_{\mathbf{N}}\) 的交，其中 \(\mathbf{N}\) 遍历 M 中在 g 下稳定的子空间。由于 s（分别为 n）的半单（分别为幂零）分量 \(x \in \mathfrak{gl}(M)\) 的半单（分别为幂零）部分，且 ad s（分别为 ad n）是 ad x 的半单（分别为幂零）部分（\(\S\) 5，第 4 号，引理 2），显然 \(x \in g^{*}\) 蕴含 \(s \in g^{*}\) 和 \(n \in g^{*}\)；因此只需证明 \(g^{*} = g\)。由于 g 是半单理想，\(g^{*}, g^{*} = a \times g\)（第 1 号，命题 1 的推论 1）。令 \(a \in a\)，令 N 为 M 中在 g 下稳定的非零子空间中极小的一个。根据 Burnside 定理，a 在 N 上的限制是恒等变换的数乘倍；按构造它的迹为零，又因 K 的特征为 0，故该限制为零。由于 M 是这类 N 的子空间的直和，故 a = 0，从而 \(g^{*} = g\)。

推论。元素 \(x\) 属于 \(\mathfrak{g}\)，它是 \(\mathbf{M}\) 的半单（分别为幂零）自同态，当且仅当 \(\operatorname{ad}_{\mathfrak{g}} x\) 是 \(\mathfrak{g}\) 的半单（分别为幂零）自同态。

令 \(s\)（分别为 \(n\)）为 \(x \in \mathfrak{g}\) 的半单（分别为幂零）分量。则 \(s \in \mathfrak{g}\) 且 \(n \in \mathfrak{g}\)（命题 3）。于是 \(\mathrm{ad}_{\mathfrak{g}} s\)（分别为 \(\mathrm{ad}_{\mathfrak{g}} n\)）是 \(\mathrm{ad}_{\mathfrak{g}} x\) 的半单（分别为幂零）分量，根据§5，第 4 号的引理 2。若 \(x\) 是半单的（分别为幂零的），则 \(\mathrm{ad}_{\mathfrak{g}} x\) 也如此。现在若 \(\mathrm{ad}_{\mathfrak{g}} x\) 是半单的（分别为幂零的），它等于 \(\mathrm{ad}_{\mathfrak{g}} s\)（分别为 \(\mathrm{ad}_{\mathfrak{g}} n\)），从而 \(x = s\)（分别为 \(x = n\)），因为 \(\mathfrak{g}\) 的伴随表示是忠实的。

定义 3. 设 g 是半单李代数。若对于每个 K 上的有限维 g-模 M，\(x_{M}\) 都是 M 的半单（分别为幂零）自同态，则称 g 的元素 x 是半单的（分别为幂零的）。

命题 4. 设 \(\mathfrak{g}\)、\(\mathfrak{g}'\) 是半单李代数，\(f\) 是从 \(\mathfrak{g}\) 到 \(\mathfrak{g}'\) 的同态。若 \(x \in \mathfrak{g}\) 是半单的（分别为幂零的），则 \(f(x)\) 也是如此。若 \(f\) 是满射，则 \(\mathfrak{g}'\) 的每个半单（分别为幂零）元素都是在 \(f\) 下由 \(\mathfrak{g}\) 中某个半单（分别为幂零）元素映出的像。

若 \(\rho\) 是 \(\mathfrak{g}'\) 的表示，则 \(\rho \circ f\) 是 \(\mathfrak{g}\) 的表示，这就得到第一个断言。若 \(f\) 是满射，则存在一个同态 \(g\)，它把 \(\mathfrak{g}'\) 映入 \(\mathfrak{g}\)，使得 \(f \circ g\) 是 \(\mathfrak{g}'\) 的恒等同态（第 1 号，命题 1 的推论 2），于是第二个断言由第一个断言得到。

定理 3. 设 g 是半单李代数。

\begin{enumerate}
\def\labelenumi{(\alph{enumi})}
\item
  设 \(x \in \mathfrak{g}\)。若存在 \(\rho\) 的忠实表示 \(\mathfrak{g}\)，使得 \(\rho(x)\) 是半单（分别为幂零）自同态，则 \(x\) 是半单的（分别为幂零的）。
\item
  \(\mathfrak{g}\) 的每个元素都能唯一地写成一个半单元素和一个幂零元素之和，且二者彼此可交换。
\end{enumerate}

设 (a) 的假设成立。令 \(\sigma\) 为 \(\mathfrak{g}\) 的一个表示，\(\mathfrak{b}\) 为 \(\sigma\) 的核的补理想，\(\alpha\) 为 \(\mathfrak{g}\) 到 \(\mathfrak{b}\) 上的投影。由命题 3 的推论，\(\mathrm{ad}_{\mathfrak{g}}x\) 是半单的（分别为幂零的），从而 \(\mathrm{ad}_{\mathfrak{b}}\alpha(x)\) 是半单的（分别为幂零的）。由于 \(\sigma(x)=\sigma(\alpha(x))\)，第一个断言由命题 3 的推论得到。第二个结果则由命题 3 应用于一个忠实表示得到。

